% !TEX root = ../main.tex
\section{Background and theory}

\hypertarget{dynamic-inconsistency-and-bellmans-principle}{%
\subsection{Dynamic inconsistency and Bellman's principle}\label{dynamic-inconsistency-and-bellmans-principle}}

Consider a planner who at time zero chooses a sequence of decisions -- here, harvest levels by stratum (defined, e.g., by species, management prescription, and age) and period -- over a finite horizon of \(T\) periods, so as to maximize an objective (here, net present value, NPV) subject to constraints, including a harvest-flow constraint that links consecutive periods. The open-loop solution is a complete plan: decisions for every period, chosen once by today's decision maker at time zero from the initial forest state.

Now suppose the plan is executed one period at a time by a sequence of planners, each inheriting the forest state realized by her predecessors and each solving the \emph{same} optimization problem from that state over the remaining horizon. Bellman's principle of optimality states that an optimal policy has the property that, whatever the initial state and initial decision, the remaining decisions must constitute an optimal policy with regard to the state resulting from the first decision \citep{bellman1957dynamic}. A plan that satisfies this property is \emph{dynamically consistent}: its tail is optimal for every subproblem it passes through. Each successive planner, re-solving from the realized state, voluntarily continues the plan. A plan that fails the property is \emph{dynamically inconsistent}: at some point the re-solving planner finds the plan's tail suboptimal for the state she has inherited and deviates from it. The period-by-period harvest trajectory of the open-loop solution is, therefore, not a forecast of realized harvests, but a projection under the counterfactual assumption of full precommitment to the time-zero plan.

Formally, let \(x = (x_1, \ldots, x_T)\) be a sequence of decisions, and let the open-loop problem from state \(s_\tau\) over the remaining horizon be \(V_\tau(s_\tau) = \max_{x_\tau, \ldots, x_T} \sum_{t=\tau}^{T} \delta^{\,t-\tau} \, u_t(x_t)\) subject to the model's constraints, where \(\delta\) is the per-period discount factor and $ u_t(x_t)$ denotes the corresponding period-$t$ contribution to the objective. The open-loop plan announced at time zero is \(x^{0} = (x_1^{0}, \ldots, x_T^{0}) = \arg\max V_1(s_1)\). Under sequential replanning, the planner at period \(\tau\) observes the realized state \(s_\tau = f(s_{\tau-1}, x_{\tau-1})\) (the forest left by her predecessor's decision) and re-solves the maximization problem from that state, \(V_\tau(s_\tau)\). The plan \(x^{0}\) is dynamically consistent if its tail is optimal for every subproblem it passes through, i.e., if \((x_\tau^{0}, \ldots, x_T^{0})\) remains a maximizer of \(V_\tau(s_\tau)\) for all \(\tau\); otherwise, some planner deviates and the plan is not followed.

These considerations have practical implications, as long-term forest plans are used to set allowable cuts, to justify regulation of harvest flow, and to price constraints and land uses through shadow prices. If the plan is dynamically inconsistent, the projected harvest levels will not be realized, the projected residual forest will not exist, and the shadow prices will characterize trade-offs for an optimal solution that will not be followed, limiting their validity as measures of trade-offs along the realized trajectory.

\hypertarget{the-open-loop-lp-forest-planning-problem}{%
\subsection{The open-loop LP forest-planning problem}\label{the-open-loop-lp-forest-planning-problem}}

The standard strategic timber harvest-scheduling model is strata-based. The forest is partitioned into development types (strata) defined by attributes such as site, species composition, age, and management prescription. The decision variables allocate stratum area to management activities over a horizon of discrete periods. The classical formulations are those of \citet{johnson1977techniques}. In \emph{Model~I}, each existing stratum is assigned a complete management schedule as a distinct activity, so that the trajectory of every initial cohort is represented explicitly. In \emph{Model~II}, decisions are expressed more compactly as area flows into a smaller set of prescription activities; \emph{Model~III} aggregates further into a network flow over age classes. We refer the reader to \citet[ch.~16]{gunn2007models} for a careful statement of the Model~I/II/III distinction.

A word on model form and size is warranted because a misconception about it circulates. Model~II is widely described as the more compact formulation due to its smaller variable count (Johnson and Scheurman 1977), and it is sometimes asserted that Model~II reduces the size and difficulty of harvest-scheduling LPs relative to Model~I. That expectation does not hold for some modern instances. In the first published comparison of the two formulations on realistic spatial problems, \citet{martin2017comparing} showed that Model~II's apparent compactness advantage may even reverse once the problem involves realistic spatial resolution or many strata. Model~II's per-strata constraint network makes its matrix denser. Its solution time grows rapidly with added structure, while Model~I's solution time grows only marginally. On their full case study, Model~I reached comparable optimal solutions in roughly a tenth of Model~II's solution time. The authors conclude that the common expectation that Model~II is the smaller, easier model is refuted for all but the smallest or least structured problems. This bears directly on the reproduction's formulation choice in our study.

\citet{daugherty1991credibility} conducted his analysis using a Model~II formulation. The reproduction reported here presents the case study as a Model~I model in \texttt{ws3}. The distinction matters for computational detail but not for the consistency question, as we demonstrate through our reproduction and scenario analysis. In our study, the Model~I choice yields the described computational advantages in realistic problem sizes. We nonetheless note the formulation difference explicitly as a matter of fidelity to the original study rather than a substantive deviation.

Two design decisions in the standard formulation are important for this analysis. First, the objective is NPV: timber values and costs in each period are weighted by a discount factor, so early harvest income is preferred to late harvest income, all else being equal. Second, harvest-flow constraints, typically even-flow or bounded-deviation constraints on total volume between consecutive periods, create a direct link between decisions in different periods. The economics of such harvest-flow constraints -- the even-flow constraint and the allowable-cut effect it enables -- have a long critical literature \citep{thompson1966traditional,schweitzer1972allowable,luckert1995allowable,hegan2000economic}. Without such a between-period link, the multi-period problem decomposes, and open-loop and sequential solutions coincide. With it, the current period's decision changes the feasible set of every future subproblem.

Writing \(H_t\) for the total harvest volume in period \(t\), without distinguishing among strata, the open-loop problem maximizes discounted net revenue subject to a harvest-flow constraint linking consecutive periods:
\begin{equation}
\max_{H_1, \ldots, H_T \geq 0} \;\; \sum_{t=1}^{T} \delta^{\,t} \, r_t \, H_t
\qquad \text{s.t.} \qquad
(1 - \alpha)\, H_t \;\le\; H_{t+1} \;\le\; (1 + \beta)\, H_t, \;\; t = 1, \ldots, T-1,
\label{eq:openloop}
\end{equation}
where \(r_t\) is net revenue per unit volume in period \(t\) and \(\delta = (1+i)^{-1}\) is calculated from the interest rate $i$. The harvest-flow constraint is the classic FORPLAN bounded-deviation ("sequential flow") form \citep[eq.~3-4, 3-5]{daugherty1991credibility}: \(\alpha\) is the maximum allowed fractional \emph{decrease} and \(\beta\) the maximum allowed fractional \emph{increase} between consecutive periods. Its special cases span the standard harvest-flow policies: non-declining yield (\(\alpha = 0\), no increase bound; the yield may not fall), bounded decline (\(\alpha > 0\), no increase bound), and symmetric bounded deviation (\(\alpha = \beta > 0\)). This between-period link prevents the multi-period problem from decomposing into independent single-period problems.

\hypertarget{daughertys-thesis}{%
\subsection{Dynamic inconsistency of open-loop forest plans}\label{daughertys-thesis}}

\citet{daugherty1991credibility} posed the credibility question directly: are the plans produced by LP-based forest-planning models, solved open-loop, plans that a sequence of future planners operating the same model would actually follow? The thesis (i) formalized dynamic consistency for LP forest-planning models, (ii) derived conditions under which open-loop optimal plans fail the consistency requirements, and (iii) demonstrated the phenomenon numerically by iterative LP simulation of sequential replanning on a real case study -- the Umpqua National Forest FORPLAN planning model documented in \citep{usda1987umpqua} -- across a designed experiment varying the initial forest condition, the harvest policy, and the interest rate. The thesis found dynamic inconsistency to be pervasive rather than exceptional, and identified the combination of the harvest-flow constraint with a disequilibrium initial forest containing financially over-mature and negatively valued strata as the structural source of the inconsistency. \citet[pp.~331--332]{gunn2007models}, citing \citet{daugherty1991credibility}, summarizes the practical symptom memorably as the ``declining non-declining yield'': a plan that promises a non-declining harvest trajectory while the sequential re-optimization it invites systematically revises the attainable harvest downward.
